\begin{mistake}{2026-04-12}{27}

\begin{problemcontent}
设微分方程 \( y''+by'+cy=0 \) 中 \( b>0,\ c>0 \)。证明其任意解 \( y(x) \) 都满足 \( \lim_{x\to +\infty} y(x) \) 存在，并求该极限。
\end{problemcontent}

\begin{solutioncontent}
对应特征方程为
\[
r^2+br+c=0
\]
设其根为 \(r_1,r_2\)，则
\[
r_1+r_2=-b<0,\qquad r_1r_2=c>0
\]
分三种情形讨论。

1. 当 \(\Delta=b^2-4c>0\) 时，方程有两个不同实根。由 \(r_1r_2>0\) 知两根同号，再由 \(r_1+r_2<0\) 知两根都小于 0。通解为
\[
y=C_1e^{r_1x}+C_2e^{r_2x}
\]
由于 \(r_1<0,r_2<0\)，故当 \(x\to +\infty\) 时，\(e^{r_1x}\to 0,\ e^{r_2x}\to 0\)，于是 \(y(x)\to 0\)。

2. 当 \(\Delta=0\) 时，方程有二重实根 \(r=-\dfrac b2<0\)。通解为
\[
y=(C_1+C_2x)e^{rx}
\]
因为 \(r<0\)，指数因子 \(e^{rx}\to 0\)，且它的衰减快于多项式 \(x\) 的增长，所以 \((C_1+C_2x)e^{rx}\to 0\)。故 \(y(x)\to 0\)。

3. 当 \(\Delta<0\) 时，方程有共轭复根
\[
r_{1,2}=-\frac b2\pm i\frac{\sqrt{4c-b^2}}2
\]
通解可写为
\[
y=e^{-\frac b2x}\left(C_1\cos \frac{\sqrt{4c-b^2}}2x+C_2\sin \frac{\sqrt{4c-b^2}}2x\right)
\]
其中三角函数有界，而 \(-\dfrac b2<0\)，故 \(e^{-\frac b2x}\to 0\)，从而 \(y(x)\to 0\)。

综上，无论特征根属于哪一类，任意解都满足
\[
\lim_{x\to +\infty}y(x)=0
\]
因此该极限存在，且极限值为 \(0\)。
\end{solutioncontent}

\mistakeknowledge{
\begin{itemize}
  \item \textbf{核心公式}：二阶常系数齐次方程先写特征方程 \(r^2+br+c=0\)；通解形式由特征根类型决定。
  \item \textbf{易错点}：由 \(r_1+r_2=-b<0\)、\(r_1r_2=c>0\) 应判断两实根都为负，不能只看到 \(b>0\) 就仓促下结论；重根情形中 \(xe^{rx}\to 0\) 也要说明。
  \item \textbf{关联知识}：判断线性常系数方程解的无穷远性态，本质上看指数项 \(e^{\lambda x}\) 中特征根实部 \(\operatorname{Re}\lambda\) 的符号。
\end{itemize}}

\end{mistake}
